% !TEX root = ../main.tex

\chapter{Méthodologie et environnement de travail}

\section{Méthodologie de travail}

Le projet est mené selon la méthodologie Scrum, avec des sprints d'une
durée de deux semaines. Cette organisation rythme l'ensemble du travail
de l'équipe et structure la planification des développements.

Chaque journée est ponctuée de deux daily meetings :

\begin{itemize}
  \item un premier avec le bureau ISI.nc, permettant de faire le point
        sur l'avancement individuel et les éventuels blocages internes ;
  \item un second avec l'équipe FPG (Product Owner, développeurs, tech
        lead), pour synchroniser l'avancement du sprint côté client.
\end{itemize}

En début de sprint, des séances de \textbf{Poker Planning} sont organisées
pour estimer collectivement la complexité des tickets à traiter, favorisant
une meilleure répartition de la charge de travail au sein de l'équipe.

Le suivi des tickets est assuré via l'outil \textbf{OpenProject}
\imgw{logo_open_project.jpeg}{1.5cm}, qui centralise le backlog, les
sprints en cours et l'état d'avancement de chaque tâche. La communication
interne quotidienne, quant à elle, s'effectue via \textbf{Mattermost}
\imgw{logo_mattermost.png}{1.5cm}.

Les fonctionnalités sont décomposées en \textit{User Stories}, détaillant
les besoins fonctionnels du point de vue de l'utilisateur final. Ce
découpage a par exemple été appliqué à la fonctionnalité d'historisation,
cœur de ma mission.

\subsection{Organisation Git dédiée à la fonctionnalité d'historisation}
\label{sec:git-workflow-historisation}

Compte tenu de l'ampleur de la fonctionnalité d'historisation, qui
touche successivement plusieurs entités et plusieurs onglets du logiciel,
une organisation Git spécifique a été mise en place afin de ne pas
perturber les autres développements en cours sur le projet.

Une branche dédiée, nommée \texttt{history}, a été créée pour centraliser
l'ensemble des développements liés à cette fonctionnalité. Un
environnement de recette en ligne a également été mis à disposition pour
cette branche, accessible à l'adresse suivante~:

\begin{center}
  \url{https://gizmo-frontend.fpg-history.dev.isi.nc/collaborateur/demandes-de-financement}
\end{center}

Pour chaque onglet ou entité à historiser, une branche fille dédiée était
créée à partir de la branche \texttt{history}, suivant la convention de
nommage \texttt{feat/nom\_de\_l'onglet\_history} — par exemple,
\texttt{feat/montage\_financier\_history} pour l'historisation de
l'onglet Montage financier.

\subsection{Méthode de travail personnelle par User Story}

Chaque User Story correspondant à l'historisation d'un onglet a été
traitée selon une méthode de travail que j'ai progressivement stabilisée
au fil du stage~:

\begin{enumerate}
  \item établir la liste des problèmes rencontrés et des écarts entre
        l'implémentation en cours et les critères d'acceptation de la
        User Story~;
  \item corriger ces problèmes un par un jusqu'à ce que l'ensemble des
        critères soit satisfait~;
  \item une fois l'implémentation jugée complète, retester
        \textbf{l'intégralité} du parcours fonctionnel concerné, de bout
        en bout~;
  \item exécuter la commande \texttt{pre-commit run --all-files} afin de
        vérifier la conformité de la syntaxe et du formatage du code
        avant tout commit~;
  \item commiter et pousser les changements sur la branche dédiée créée
        en amont~;
  \item ouvrir une \textit{merge request} de cette branche vers la
        branche \texttt{history}~;
  \item soumettre cette \textit{merge request} à la vérification de mon
        tuteur~;
  \item une fois validée, fusionner les changements sur la branche
        \texttt{history} et sur l'environnement de recette associé~;
  \item retester la fonctionnalité directement sur cet environnement de
        recette~;
  \item si l'ensemble des tests est concluant, passer à la User Story
        suivante — chaque User Story correspondant à l'historisation d'un
        onglet distinct.
\end{enumerate}

Cette méthode, bien que rigoureuse et parfois répétitive, a permis de
limiter les régressions entre les différentes User Stories traitées
successivement, chacune venant s'appuyer sur un état stable et validé de
la branche \texttt{history}.

\section{Stack technique}

\subsection{Développement}
\begin{itemize}
  \item \textbf{VS Code} \imgw{logo_vs_code.png}{1.2cm}  -  IDE principal
        utilisé pour l'ensemble des développements, frontend comme backend.
\end{itemize}

\subsection{Frontend}
\begin{itemize}
  \item \textbf{Angular} \imgw{logo_angular.jpeg}{1.2cm}  -  framework
        principal du frontend, associé à \textbf{PrimeNG} comme
        librairie de composants UI, offrant des composants riches
        (tableaux, formulaires, timeline, etc.) directement exploitables.
  \item \textbf{Gizmo} \imgw{logo_gizmo.png}{1.2cm}  -  application
        constituant l'interface destinée aux structures adhérentes.
\end{itemize}

\subsection{Backend}
\begin{itemize}
  \item \textbf{Python / Django} \imgw{logo_django_python.png}{1.5cm},
        associé à \textbf{Django REST Framework} pour l'exposition d'une
        API REST consommée par le frontend Angular.
  \item \textbf{django-simple-history}  -  bibliothèque permettant
        l'historisation automatique des modèles Django, point central
        de la mission qui m'a été confiée et détaillée au
        Chapitre~\ref{chap:historisation}.
\end{itemize}

\subsection{Base de données}
\begin{itemize}
  \item \textbf{PostgreSQL} \imgw{logo_postgres.png}{1.2cm}  -  système de
        gestion de base de données relationnelle utilisé pour l'ensemble
        du projet.
  \item \textbf{DBeaver} \imgw{logo_dbeaver.jpeg}{1.2cm}  -  client de base
        de données utilisé pour l'inspection et le débogage direct des
        tables, notamment des tables d'historique.
\end{itemize}

\subsection{DevOps et infrastructure}
\begin{itemize}
  \item \textbf{GitLab} \imgw{logo_gitlab.png}{1.2cm}  -  versionning du
        code, intégration et déploiement continus (CI/CD), gestion des
        \textit{merge requests} et revues de code.
  \item \textbf{Linux} \imgw{logo_linux.png}{1.2cm}  -  système
        d'exploitation utilisé comme environnement de développement.
  \item \textbf{Docker} \imgw{logo_docker.png}{1.2cm}  -  conteneurisation
        des différents services du projet : base de données, stockage de
        fichiers (MinIO), backend Django.
\end{itemize}

\subsection{Outils découverts pendant le stage}

Au-delà des outils déjà maîtrisés en amont du stage, ce dernier a été
l'occasion de découvrir de nouveaux outils, venus enrichir mon confort de
travail au quotidien.

\begin{itemize}
  \item \textbf{LazyGit} \imgw{logo_lazy_git.jpeg}{1.2cm}  -  interface en
        ligne de commande offrant une vision claire et interactive de
        l'état d'un dépôt Git (branches, historique des commits,
        différences entre fichiers). Cet outil a nettement facilité la
        gestion quotidienne des nombreuses branches créées dans le cadre
        du workflow Git dédié à l'historisation, détaillé en
        Section~\ref{sec:git-workflow-historisation}.
\end{itemize}

\section{Architecture générale du projet}

L'architecture du projet suit un modèle classique d'application web
découplée :

\begin{itemize}
  \item un \textbf{frontend Angular} qui consomme une \textbf{API REST
        Django} exposée via Django REST Framework ;
  \item un \textbf{stockage de fichiers} (documents justificatifs,
        pièces jointes) assuré par \textbf{MinIO}, une solution de
        stockage objet compatible avec le protocole S3 ;
  \item une \textbf{génération de rapports PDF} déléguée à un serveur
        \textbf{JasperReports} externe, sollicité par le backend pour la
        production de documents officiels (attestations, décisions de
        financement, etc.).
\end{itemize}

\begin{figure}[H]
  \centering
  \imgw{schema_architecture.png}{0.8\textwidth}
  \caption{Schéma de l'architecture générale du projet}
  \label{fig:architecture}
\end{figure}

Cette architecture découplée permet une séparation claire des
responsabilités : le frontend se concentre sur l'expérience utilisateur,
le backend sur la logique métier et l'accès aux données, tandis que les
services externes (MinIO, JasperReports) prennent en charge des besoins
spécialisés sans alourdir le cœur applicatif.
